\documentclass[11pt,a4paper]{article}

% Template visual Matriz Aprova. Compile com XeLaTeX ou LuaLaTeX.
\usepackage{fontspec}
\usepackage{polyglossia}
\setmainlanguage{portuguese}
\IfFontExistsTF{Aptos}{
  \setmainfont{Aptos}
  \setsansfont{Aptos}
}{
  \IfFontExistsTF{Calibri}{
    \setmainfont{Calibri}
    \setsansfont{Calibri}
  }{
    \IfFontExistsTF{TeX Gyre Heros}{
      \setmainfont{TeX Gyre Heros}
      \setsansfont{TeX Gyre Heros}
    }{
      \setmainfont{Latin Modern Sans}
      \setsansfont{Latin Modern Sans}
    }
  }
}
\usepackage[a4paper,margin=2cm,headheight=16pt]{geometry}
\usepackage{graphicx}
\usepackage{xcolor}
\usepackage{tikz}
\usepackage{tcolorbox}
\usepackage{enumitem}
\usepackage{booktabs}
\usepackage{tabularx}
\usepackage{amssymb}
\usepackage{fancyhdr}
\usepackage{titlesec}
\usepackage{hyperref}

\definecolor{matrizlime}{HTML}{CBFF4D}
\definecolor{matrizink}{HTML}{101313}
\definecolor{matrizmuted}{HTML}{596363}
\definecolor{matrizline}{HTML}{DDE3DF}
\definecolor{matrizsoft}{HTML}{F3F7EC}

\hypersetup{colorlinks=true,linkcolor=matrizink,urlcolor=matrizink}
\pagestyle{fancy}
\fancyhf{}
\lhead{\textsf{\small MATRIZ APROVA}}
\rhead{\textsf{\small APOSTILA}}
\cfoot{\textsf{\small \thepage}}
\renewcommand{\headrulewidth}{0.4pt}
\renewcommand{\headrule}{\hbox to\headwidth{\color{matrizline}\leaders\hrule height \headrulewidth\hfill}}

\titleformat{\section}{\Large\bfseries\color{matrizink}}{\thesection}{0.6em}{}
\titleformat{\subsection}{\large\bfseries\color{matrizink}}{\thesubsection}{0.6em}{}
\setlist{nosep,leftmargin=*}

\newtcolorbox{foco}{colback=matrizsoft,colframe=matrizlime,boxrule=1pt,arc=3pt,left=10pt,right=10pt,top=8pt,bottom=8pt}
\newtcolorbox{lei}{colback=matrizink,colframe=matrizink,coltext=white,boxrule=0pt,arc=3pt,left=10pt,right=10pt,top=8pt,bottom=8pt}
\newtcolorbox{alerta}{colback=white,colframe=matrizline,boxrule=0.6pt,arc=3pt,left=10pt,right=10pt,top=8pt,bottom=8pt}

% Logo usada na capa. Caminho relativo à pasta onde o .tex é compilado
% (materiais/<disciplina>/). Ajuste se a estrutura mudar.
\newcommand{\logomatriz}{../../public/matrizaprova_temaclaro.png}

\begin{document}

\begin{titlepage}
  \thispagestyle{empty}
  \begin{center}
    \includegraphics[width=0.55\linewidth]{\logomatriz}\par
  \end{center}
  \vspace{2.2cm}
  {\sffamily\fontsize{28}{32}\selectfont\bfseries Apostila — Física: Cinemática e dinâmica\par}
  \vspace{0.4cm}
  {\sffamily\large ESA, EsPCEx, EEAR, AFA, EFOMM, Escola Naval e ENEM\par}
  \vfill
  \begin{foco}
    \textsf{\textbf{Foco da apostila}}\\[3pt]
    escolha um eixo, indique o sentido positivo e trabalhe com grandezas compatíveis. Um sinal negativo informa sentido ou desaceleração; não é um erro automático.
  \end{foco}
  \vspace{0.7cm}
  {\sffamily\small Atualizado em 16 de agosto de 2026\quad ·\quad adaptação necessária ao edital e ao nível da prova\par}
  \vspace{1cm}
  {\sffamily\small matrizaprova.com\par}
\end{titlepage}

\tableofcontents
\newpage

% Conteúdo gerado entra aqui.
\section*{O que cai}

Cinemática descreve o movimento. Dinâmica explica suas causas por meio das forças. Em provas militares e no ENEM, os cálculos costumam ser diretos, mas o maior risco está em escolher a equação errada, misturar unidades ou ignorar o sentido dos vetores.

Os pontos de maior incidência são:

\begin{enumerate}
  \item conversão de unidades e leitura de gráficos de posição, velocidade e tempo;
  \item movimento uniforme e movimento uniformemente variado;
  \item queda livre e lançamento vertical;
  \item leis de Newton, força resultante e diagramas de corpo livre;
  \item atrito, peso, normal e aplicação da segunda lei em planos horizontais e
\end{enumerate}

   inclinados.

\textbf{Regra de prova:} escolha um eixo, indique o sentido positivo e trabalhe com grandezas compatíveis. Um sinal negativo informa sentido ou desaceleração; não é um erro automático.

\section*{Teoria essencial}

\subsection*{1. Referencial, posição e deslocamento}

Movimento e repouso dependem do referencial. Uma pessoa sentada em um ônibus está em repouso em relação ao ônibus, mas em movimento em relação à rua.

A posição é indicada por uma coordenada, geralmente s. O deslocamento é a variação da posição:

deslocamento = posição final − posição inicial.

O deslocamento é vetorial e pode ser positivo, negativo ou nulo. A distância percorrida é escalar e corresponde ao comprimento total do caminho. Se um atleta completa uma volta e retorna ao ponto de partida, o deslocamento é zero, mas a distância percorrida não é.

A velocidade média é o deslocamento dividido pelo intervalo de tempo. Já a rapidez média usa a distância total. Em uma ida e volta, a velocidade média pode ser zero, enquanto a rapidez média permanece positiva.

\subsection*{2. Unidades e conversões}

No Sistema Internacional, posição é medida em metro, tempo em segundo, velocidade em metro por segundo e aceleração em metro por segundo ao quadrado.

As conversões mais cobradas são:

\begin{itemize}
  \item de km/h para m/s: divida por 3,6;
  \item de m/s para km/h: multiplique por 3,6;
  \item de minuto para segundo: multiplique por 60;
  \item de hora para segundo: multiplique por 3.600.
\end{itemize}

Exemplo: 72 km/h correspondem a 72 ÷ 3,6 = 20 m/s. Não se deve substituir diretamente 72 na equação que usa metro e segundo.

\subsection*{3. Movimento uniforme}

No movimento uniforme, a velocidade é constante e a aceleração é nula. A posição varia de forma linear com o tempo:

posição final = posição inicial + velocidade × tempo.

Se a posição inicial for zero, basta multiplicar velocidade por tempo. O gráfico posição × tempo é uma reta, e sua inclinação representa a velocidade. Reta crescente indica velocidade positiva; reta decrescente indica velocidade negativa; reta horizontal indica repouso.

Em problemas de encontro, escreva uma posição para cada móvel usando o mesmo referencial. O encontro ocorre quando as duas posições são iguais. Evite usar apenas “velocidade relativa” quando os sentidos não estiverem claros.

\subsection*{4. Movimento uniformemente variado}

No movimento uniformemente variado, a aceleração é constante. As relações mais usadas são:

\begin{itemize}
  \item velocidade final = velocidade inicial + aceleração × tempo;
  \item posição final = posição inicial + velocidade inicial × tempo + metade da aceleração × tempo ao quadrado;
  \item velocidade final ao quadrado = velocidade inicial ao quadrado + duas vezes a aceleração × deslocamento.
\end{itemize}

A terceira relação, conhecida como equação de Torricelli, é útil quando o tempo não aparece no enunciado.

Se velocidade e aceleração têm o mesmo sinal, o módulo da velocidade aumenta. Se têm sinais opostos, o móvel desacelera naquele instante. A aceleração negativa não significa sempre que o móvel está diminuindo a rapidez: isso depende do sinal da velocidade.

No gráfico velocidade × tempo, a inclinação representa a aceleração e a área algébrica representa o deslocamento. No gráfico aceleração × tempo, a área representa a variação da velocidade.

\subsection*{5. Queda livre e lançamento vertical}

Desprezando a resistência do ar, todos os corpos próximos à superfície da Terra possuem a mesma aceleração gravitacional, aproximadamente 10 m/s² em muitos exercícios. O valor exato indicado no enunciado deve prevalecer.

Na queda livre, o corpo parte do repouso e acelera no sentido da gravidade. Se o eixo vertical positivo for para cima, a aceleração é negativa. No lançamento vertical para cima, a velocidade diminui até chegar a zero no ponto mais alto. Depois, o corpo inicia a descida com velocidade crescente.

No ponto mais alto:

\begin{itemize}
  \item a velocidade instantânea é zero;
  \item a aceleração continua sendo a gravidade;
  \item o corpo não fica sem força gravitacional.
\end{itemize}

Quando o lançamento e a queda ocorrem no mesmo nível e não há resistência do ar, o tempo de subida é igual ao tempo de descida, e o módulo da velocidade de retorno é igual ao módulo da velocidade inicial.

\subsection*{6. Vetores e força resultante}

Força é uma grandeza vetorial: possui módulo, direção e sentido. A força resultante é a soma vetorial de todas as forças aplicadas ao corpo.

Em uma dimensão, escolha um sentido positivo e use sinais. Duas forças de mesmo módulo e sentidos opostos têm resultante nula. Em direções perpendiculares, a resultante pode ser calculada pelo teorema de Pitágoras quando as forças são componentes de um mesmo vetor.

O diagrama de corpo livre deve representar apenas as forças que atuam no corpo analisado. Não desenhe a força que o corpo exerce em outro objeto como se ela atuasse nele.

\subsection*{7. Leis de Newton}

\subsubsection*{Primeira lei: inércia}

Se a força resultante é nula, o corpo permanece em repouso ou em movimento retilíneo uniforme. Força resultante nula não significa ausência de forças; forças podem existir e se equilibrar.

\subsubsection*{Segunda lei: princípio fundamental}

A força resultante é igual ao produto da massa pela aceleração:

força resultante = massa × aceleração.

No Sistema Internacional, força é medida em newton, equivalente a quilograma vez metro por segundo ao quadrado. A mesma força produz maior aceleração em um corpo de menor massa.

\subsubsection*{Terceira lei: ação e reação}

As forças de ação e reação têm mesmo módulo, mesma direção e sentidos opostos, mas atuam em corpos diferentes. Por isso, não se anulam no diagrama de um único corpo. O pé empurra o chão; o chão empurra o pé com força de mesmo módulo.

\subsection*{8. Peso, normal e tração}

O peso é a força gravitacional exercida pela Terra:

peso = massa × gravidade.

Peso e massa são grandezas diferentes. A massa é uma propriedade do corpo e não depende do local; o peso depende do campo gravitacional.

A força normal é perpendicular à superfície de contato. Em uma superfície horizontal, com ausência de aceleração vertical e sem outras forças verticais, seu módulo é igual ao peso. Isso não vale automaticamente em um plano inclinado ou quando há outras forças verticais.

A tração é uma força transmitida por fio ou cabo ideal. Em um fio sem massa e inextensível, a tração tem o mesmo módulo ao longo do fio, mas corpos ligados podem ter acelerações diferentes se houver polias ou restrições distintas.

\subsection*{9. Atrito e plano inclinado}

O atrito se opõe ao movimento ou à tendência de movimento relativo entre as superfícies. O atrito estático atua antes do deslizamento e pode assumir valores até um limite máximo. O atrito cinético atua durante o deslizamento.

Em um modelo usual:

\begin{itemize}
  \item atrito estático máximo = coeficiente estático × normal;
  \item atrito cinético = coeficiente cinético × normal.
\end{itemize}

O atrito não é sempre igual ao coeficiente multiplicado pela normal: essa igualdade vale para o atrito cinético e para o valor máximo do estático. Se o corpo está parado e a força horizontal aplicada é pequena, o atrito estático se ajusta para equilibrá-la.

No plano inclinado, decomponha o peso em duas componentes: uma perpendicular ao plano, relacionada à normal, e outra paralela ao plano, responsável pela tendência de deslizamento. O ângulo e o eixo escolhido devem ser desenhados antes do cálculo.

\section*{Como a banca cobra}

\begin{itemize}
  \item Fornece velocidade em km/h e pede aceleração em m/s² sem destacar a conversão.
  \item Mostra um gráfico e troca inclinação por área: no gráfico s × t, inclinação é velocidade; no v × t, área é deslocamento.
  \item Afirma que aceleração zero significa repouso, confundindo com movimento uniforme.
  \item Diz que no ponto mais alto do lançamento vertical a aceleração é zero.
  \item Faz a força de ação e reação atuar no mesmo corpo para induzir erro.
  \item Considera normal igual ao peso em qualquer situação, inclusive em planos inclinados.
  \item Trata o atrito estático como sempre igual ao seu valor máximo.
\end{itemize}

\subsection*{Exemplos resolvidos}

\subsubsection*{Exemplo 1}

Um veículo parte do repouso e acelera constantemente a 2 m/s² durante 8 s. Determine a velocidade final e o deslocamento.

\begin{enumerate}
  \item Dados: velocidade inicial zero, aceleração 2 m/s² e tempo 8 s.
  \item Velocidade final: 0 + 2 × 8 = 16 m/s.
  \item Deslocamento: 0 + 0 × 8 + metade de 2 × 8 ao quadrado = 64 m.
\end{enumerate}

O veículo termina o intervalo a 16 m/s, depois de percorrer 64 m.

\subsubsection*{Exemplo 2}

Um bloco de massa 5 kg recebe força horizontal de 30 N para a direita e força de atrito de 10 N para a esquerda. Determine a aceleração.

\begin{enumerate}
  \item Escolha a direita como sentido positivo.
  \item Resultante: 30 − 10 = 20 N.
  \item Pela segunda lei: aceleração = 20 ÷ 5 = 4 m/s².
\end{enumerate}

O bloco acelera para a direita. O atrito reduziu a resultante, mas não impediu o movimento.

\section*{Quadro-resumo}

\begin{tabularx}{\linewidth}{llX}
\toprule
 \textbf{Situação} & \textbf{Relação principal} & \textbf{Interpretação} \\
\midrule
 movimento uniforme & posição = inicial + velocidade × tempo & velocidade constante \\
 movimento variado & velocidade final = inicial + aceleração × tempo & aceleração constante \\
 sem tempo & v² = v₀² + 2 × a × deslocamento & Torricelli \\
 gráfico v × t & área = deslocamento & área algébrica \\
 força resultante & força = massa × aceleração & segunda lei \\
 equilíbrio & resultante = zero & repouso ou velocidade constante \\
 peso & peso = massa × gravidade & força gravitacional \\
 ação e reação & mesmo módulo, corpos diferentes & não se anulam no mesmo corpo \\
\bottomrule
\end{tabularx}


\section*{Questões de fixação}

\subsection*{Questão 1 — (questão inédita)}

Um móvel desloca-se a 54 km/h com velocidade constante. Sua velocidade no Sistema Internacional é:

A) 5,4 m/s. \\ B) 15 m/s. \\ C) 19,44 m/s. \\ D) 54 m/s. \\ E) 194,4 m/s.

\begin{alerta}
\textbf{Gabarito: B.} Para converter km/h em m/s, divide-se por 3,6: 54 ÷ 3,6 = 15 m/s.
\end{alerta}

\subsection*{Questão 2 — (questão inédita)}

Um corpo está em movimento retilíneo com aceleração nula. É necessariamente correto afirmar que ele:

A) está parado. B) possui velocidade crescente. C) possui força resultante nula. D) não sofre nenhuma força. E) possui deslocamento nulo.

\begin{alerta}
\textbf{Gabarito: C.} Pela segunda lei, aceleração nula implica força resultante nula. O corpo pode estar parado ou em movimento uniforme; forças individuais podem existir e se equilibrar.
\end{alerta}

\subsection*{Questão 3 — (questão inédita)}

No ponto mais alto de um lançamento vertical, desprezando a resistência do ar, o corpo apresenta:

A) velocidade e aceleração nulas. B) velocidade máxima e aceleração nula. C) velocidade nula e aceleração igual à gravidade. D) velocidade igual à gravidade e aceleração nula. E) força resultante nula.

\begin{alerta}
\textbf{Gabarito: C.} A velocidade instantânea é zero no ponto mais alto, mas a gravidade continua atuando e produz aceleração para baixo.
\end{alerta}

\subsection*{Questão 4 — (questão inédita)}

Um bloco de 4 kg sofre força resultante horizontal de 12 N. Sua aceleração é:

A) 0,33 m/s². \\ B) 3 m/s². \\ C) 8 m/s². \\ D) 16 m/s². \\ E) 48 m/s².

\begin{alerta}
\textbf{Gabarito: B.} Pela segunda lei, aceleração = força resultante ÷ massa = 12 ÷ 4 = 3 m/s².
\end{alerta}

\subsection*{Questão 5 — (questão inédita)}

Assinale a afirmativa correta sobre ação e reação.

A) São forças que atuam sempre no mesmo corpo. B) Possuem módulos diferentes quando os corpos têm massas diferentes. C) Anulam-se porque têm sentidos opostos. D) Possuem mesmo módulo e atuam em corpos diferentes. E) Só existem quando os corpos estão em movimento.

\begin{alerta}
\textbf{Gabarito: D.} A terceira lei de Newton estabelece mesmo módulo, mesma direção e sentidos opostos, com atuação em corpos diferentes. Por isso não se anulam no diagrama de um único corpo.
\end{alerta}

\section*{Revisão final}

\begin{itemize}
  \item[$\square$] Converti todas as unidades para um sistema compatível.
  \item[$\square$] Diferenciei deslocamento de distância percorrida.
  \item[$\square$] Usei inclinação e área corretas nos gráficos.
  \item[$\square$] Escolhi o eixo e o sentido positivo antes das equações.
  \item[$\square$] Lembrei que aceleração zero não significa necessariamente repouso.
  \item[$\square$] Mantive a gravidade no ponto mais alto do lançamento.
  \item[$\square$] Desenhei somente as forças que atuam no corpo analisado.
  \item[$\square$] Separei massa, peso e força normal.
  \item[$\square$] Diferenciei atrito estático de atrito cinético.
  \item[$\square$] Conferi unidade e sinal da resposta final.
\end{itemize}

\end{document}
